
We find that a permutation-equivariant graph SSL encoder trained on CNN checkpoints achieves $R^{2}=0.231$ on held-out ViT-S/16 ImageNet models, compared to $R^{2}=0.072$ for SANE's flat tokenization --- a ~220\% relative improvement (3.$19\times$ ratio) with no ViT training data. The source of improvement is the directed computational graph schema, which provides an architecture-agnostic coordinate system that flat tokenization cannot. These results are from a single training seed with 53 test models; magnitude estimates carry uncertainty, and full statistical validation is required.

A secondary result is a principled negative: scale equivariance does not help and appears to hurt ViT cross-architecture SSL transfer ($\Delta R^{2}=-0.021$ in the controlled h-m1 comparison; $\Delta R^{2}=-0.143$ in the h-m2 ablation, partially confounded by a training duration discrepancy). We attribute this to a LayerNorm gauge-fixing hypothesis --- LayerNorm eliminates the scale degree of freedom that scale equivariance targets --- though causal confirmation requires additional experiments on non-LayerNorm target architectures.

A third finding is methodological: MMD is an unreliable cross-architecture transfer metric when baseline encoders may collapse representations. Latent collapse produces artificially low MMD without genuine alignment, making $R^{2}-based$ linear probe evaluation preferable.

\textbf{Future Directions.} (1) \textit{Normalization-conditional equivariance}: Evaluate scale vs. permutation-only equivariance on non-LayerNorm architectures to test the gauge-fixing hypothesis causally. (2) \textit{Multi-seed, full-zoo validation}: 5 seeds and 250+ ViT models for statistical significance and calibrated magnitude estimates. (3) \textit{Functional decoder design}: Hypernetwork-style decoder with architecture-aware weight output projections for latent interpolation. (4) \textit{Epoch-matched ablation}: Re-train EquiSSL for 100 epochs to provide a fully controlled symmetry group comparison.

